\section{Requisitos No Funcionales}

\subsection{RNF1 Seguridad de la Información}
Este bloque cubre los aspectos transversales de protección de la
información gestionada por el sistema: cifrado, gestión de secretos,
protección frente a vulnerabilidades comunes y compartición del
acceso.
\begin{itemize}
    \item \textbf{RNF1.1 Cifrado de la información en tránsito} \\
    Toda comunicación entre cliente y servidor debe realizarse mediante TLS 1.2 o superior.
    \begin{description}
        \item[-] El certificado TLS lo provisiona automáticamente Render mediante Let's Encrypt y se renueva sin intervención manual.
        \item[-] Las comunicaciones internas entre los servicios desplegados en
        Render (backend a base de datos) viajan por la red privada del
        proyecto, que es por sí misma una red aislada.  
    \end{description}


    \item \textbf{RNF1.2 Cifrado de datos sensibles en reposo} \\
    Los datos sensibles almacenados en la base de datos deben estar
    protegidos frente a acceso no autorizado al medio físico.
    \begin{description}
        \item[-] El cifrado a nivel de fichero del sistema lo proporciona Render
        como parte del servicio gestionado de PostgreSQL.\@
    \end{description}


    \item \textbf{RNF1.3 Gestión de secretos} \\
    Los secretos de configuración (cadena de conexión a la base de datos,
    secreto de firma JWT, claves de servicios externos cuando se añadan)
    nunca se persisten en el código fuente ni en el repositorio.
    \begin{description}
        \item[-] Los secretos se gestionan exclusivamente como variables de entorno
        cifradas en reposo en el dashboard de Render.
        \item[-] El repositorio incluye un fichero \texttt{.env.example} con los nombres de
        las variables esperadas pero sin valores reales.
    \end{description}


    \item \textbf{RNF1.4 Protección frente a vulnerabilidades web comunes} \\
    El sistema debe protegerse activamente frente a las clases de
    vulnerabilidad recogidas en la lista OWASP Top 10 vigente.
    \begin{description}
        \item[-] \textbf{Inyección SQL}: prevención mediante uso exclusivo de consultas
        parametrizadas a través de Spring Data JPA y, cuando se requieran
        consultas nativas, mediante el uso de parámetros con bind, nunca
        concatenación de cadenas. 
        \item[-] \textbf{Cross-Site Scripting (XSS)}: prevención mediante el escapado
        automático que Angular aplica por defecto a las interpolaciones
        en plantillas, y la prohibición de uso de \texttt{innerHTML} o
        \texttt{bypassSecurityTrust}.
        \item[-] \textbf{Cross-Site Request Forgery (CSRF)}: como el sistema utiliza
        autenticación basada en JWT enviado en cabecera \texttt{Authorization},
        y no en cookies de sesión, el vector clásico de CSRF queda
        mitigado por construcción.
        \item[-] \textbf{Inyección de cabeceras}: validación estricta de cabeceras
        HTTP recibidas y construcción controlada de cabeceras enviadas.
    \end{description}


    \item \textbf{RNF1.5 Compartición del acceso a datos sensibles} \\
    El acceso a datos sensibles está compartimentado por rol, nivel de
    seguridad y, cuando aplica, habilitación temporal explícita
    (RF2.3, RF2.4).
    \begin{description}
        \item[-] Las consultas a la base de datos que devuelvan datos sensibles
        filtran obligatoriamente en función del usuario autenticado: ningún
        endpoint devuelve datos sensibles sin haber verificado primero la
        autorización del solicitante.
        \item[-] Las exportaciones que incluyen datos personales de amigos
        registran el evento en el log de auditoría con identificador del
        voluntario que las realiza.
    \end{description}


    \item \textbf{RNF1.6 Análisis de dependencias} \\
    El proyecto debe revisar periódicamente sus dependencias en busca de
    vulnerabilidades conocidas.
    \begin{description}
        \item[-] Las dependencias del backend Java se revisan mediante
        \texttt{dependency-check} (OWASP) o equivalente, integrado en el pipeline
        de integración continua.
        \item[-] Las dependencias del frontend Angular se revisan mediante \texttt{npm audit}
        o equivalente, integrado en el pipeline.
        \item[-] Las vulnerabilidades de severidad alta o crítica deben corregirse
        antes del siguiente release a producción.
    \end{description}


    \item \textbf{RNF1.7 Escaneo de secretos en el repositorio} \\
    El repositorio Git debe protegerse frente a la inclusión accidental
    de secretos en el código fuente utilizando herramientas como \texttt{git-secrets}, \texttt{gitleaks} o el escaneo nativo de GitHub configurado
    como hook de pre-commit y como verificación en el pipeline.
\end{itemize}



\subsection{RNF2 Autenticación}
Este bloque cubre los mecanismos por los que el sistema verifica la
identidad de un voluntario que pretende usarlo.
\begin{itemize}
    \item \textbf{RNF2.1 Autenticación basada en credenciales} \\
    El sistema autentica a sus usuarios mediante correo electrónico + contraseña. El correo electrónico es único en el sistema (RF1.1) y actúa como
    identificador.


    \item \textbf{RNF2.2 Almacenamiento seguro de contraseñas} \\
    Las contraseñas no se almacenan en claro bajo ninguna circunstancia.
    \begin{description}
        \item[-] Se utiliza un algoritmo de derivación robusto y resistente a
        ataques por fuerza bruta: BCrypt con factor de coste configurable
        (mínimo 10). 
        \item[-] El backend nunca devuelve la contraseña almacenada ni en respuestas de API ni en logs.
        \item[-] Las contraseñas no aparecen en mensajes de error ni en notificaciones. 
    \end{description}


    \item \textbf{RNF2.3 Tokens JWT con caducidad limitada} \\
    La sesión de un voluntario autenticado se materializa en un token JWT
    con tiempo de vida limitado.
    \begin{description}
        \item[-] Tiempo de vida del token de acceso: 60 minutos por defecto, configurable. 
        \item[-] El token incluye en sus claims el identificador del voluntario, el
        rol, el nivel de seguridad y las capacidades específicas asignadas.
        \item[-] El token se firma con un secreto que reside en variables de entorno (RNF1.3).
        \item[-] No se incluyen en el token datos sensibles del usuario ni datos de amigos. 
    \end{description}


    \item \textbf{RNF2.4 Renovación de sesión} \\
    El sistema debe renovar la sesión de un voluntario activo sin obligarle
    a reintroducir credenciales en cada caducidad del token.
    \begin{description}
        \item[-] Mecanismo: token de refresco de vida más larga (por defecto, 7 días) almacenado de forma segura en el cliente. 
        \item[-] El token de refresco se invalida explícitamente al cerrar sesión (RF1.5).
        \item[-] La renovación se realiza de forma transparente para el usuario.
    \end{description}


    \item \textbf{RNF2.5 Cierre automático por inactividad} \\
    El sistema debe cerrar automáticamente la sesión de un voluntario
    tras un periodo configurable de inactividad.
    \begin{description}
        \item[-] Inactividad por defecto: 60 minutos.
        \item[-] Se considera actividad cualquier petición autenticada al backend.
        \item[-] 5 minutos antes del cierre, el frontend muestra un aviso visible.
    \end{description}
\end{itemize}


\subsection{RNF3 Autorización y Control de Acceso}
Este bloque cubre la verificación de qué puede hacer un usuario una
vez autenticado.
\begin{itemize}
    \item \textbf{RNF3.1 Modelo de control de acceso basado en roles y capacidades} \\
    El sistema implementa autorización mediante un modelo combinado de
    roles fijos (RF2.1) y capacidades granulares (RF2.2) que amplían o
    restringen el comportamiento por defecto del rol.


    \item \textbf{RNF3.2 Verificación obligatoria en backend} \\
    Toda operación que acceda a datos sensibles o modifique el estado
    del sistema debe verificar la autorización del solicitante en el
    backend.
    \begin{description}
        \item[-] La ocultación de elementos en la interfaz del frontend no
        constituye por sí sola un mecanismo de autorización. 
        \item[-] Cada endpoint expuesto declara explícitamente las capacidades
        requeridas para invocarlo, mediante anotaciones (\texttt{@PreAuthorize})
        o equivalente declarativo.
        \item[-] Una llamada sin las capacidades necesarias responde con código
        HTTP 403 sin filtrar información sobre la operación solicitada.
    \end{description}


    \item \textbf{RNF3.3 Caducidad automática de habilitaciones} \\
    Las habilitaciones temporales (RF2.4) caducan automáticamente sin intervención manual.


    \item \textbf{RNF3.4 Principio de privilegio mínimo} \\
    El diseño del sistema sigue el principio de privilegio mínimo: cada
    voluntario recibe el conjunto de capacidades estrictamente necesarias
    para desempeñar su función.


    \item \textbf{RNF3.5 Auditoría de cambios en autorización} \\
    Toda modificación del rol, nivel de seguridad o capacidades de un
    voluntario queda registrada en el log de auditoría con identificador
    del Admin que realiza el cambio, valor anterior, valor nuevo y marca
    de tiempo (véase también RNF5).
\end{itemize}


\subsection{RNF4 Validación de Datos}
Este bloque cubre la verificación de los datos introducidos en el
sistema antes de su persistencia o procesamiento.
\begin{itemize}
    \item \textbf{RNF4.1 Validación obligatoria en backend} \\
    Toda entrada de datos al sistema debe validarse en el backend con
    independencia de las validaciones que el frontend pueda realizar. La validación en el frontend mejora la experiencia de usuario pero
    nunca sustituye a la validación en el backend.


    \item \textbf{RNF4.2 Validación declarativa de DTOs} \\
    La validación de datos en el backend se implementa de forma
    declarativa mediante Bean Validation (\texttt{jakarta.validation.constraints})
    sobre los DTOs de entrada.
    \begin{description}
        \item[-] Los controladores declaran \texttt{@Valid} o \texttt{@Validated} en los
        parámetros de entrada para activar la validación automática. 
        \item[-] Las violaciones de validación se traducen en respuestas HTTP 400
        con detalle estructurado de los campos inválidos.
        \item[-] Las restricciones de validación se documentan en el código junto al DTO al que aplican.
    \end{description}


    \item \textbf{RNF4.3 Validaciones de formato} \\
    El sistema valida que los datos introducidos respetan el formato esperado.
    \begin{description}
        \item[-] Correos electrónicos cumplen la sintaxis RFC 5321/5322.
        \item[-] Fechas válidas en formato ISO 8601.
        \item[-] Importes numéricos con dos decimales máximo y separador decimal punto.   
        \item[-] Cadenas de texto con longitudes máximas explícitas.
    \end{description}


    \item \textbf{RNF4.4 Validaciones de lógica de negocio} \\
    El sistema valida no solo formato sino también las reglas de negocio
    declaradas en los requisitos funcionales.
    \begin{description}
        \item[-] Una entrega de prenda no puede registrarse si no hay stock disponible (RF5.3.3). 
        \item[-] Una salida alimentaria con prioridad distinta de 1 requiere justificación obligatoria (RF6.4). 
        \item[-] Un asiento contable con tipo Amigo requiere amigo vinculado (RF7.4.2).
        \item[-] Las reglas se implementan en la capa de servicio, no en
        controladores ni repositorios, para mantener la separación de
        responsabilidades.
    \end{description}


    \item \textbf{RNF4.5 Entradas de texto seguras y limpias en el sistema} \\
    Las entradas de texto libre destinadas a almacenarse y posteriormente
    mostrarse a otros usuarios se sanitizan para evitar inyección de
    contenido malicioso.
    \begin{description}
        \item[-] El escapado de salida en plantillas Angular ya proporciona
        protección frente a XSS por defecto. 
    \end{description}
\end{itemize}


\subsection{RNF5 Auditoría y Trazabilidad}
Este bloque cubre el registro de las operaciones que el sistema
realiza, con el fin de garantizar trazabilidad ante incidentes,
auditorías y conflictos.
\begin{itemize}
    \item \textbf{RNF5.1 Registro de operaciones sensibles} \\
    El sistema mantiene un registro persistente e inmutable de las
    operaciones que afectan a datos sensibles, datos contables o al
    modelo de autorización.
    \begin{description}
        \item[-] Operaciones registradas: creación, modificación y eliminación de
        fichas de amigo en cualquier nivel; cambio de nivel de confianza;
        edición y anulación de asientos contables; entrada y salida
        alimentaria; modificación de cuenta de voluntario; cambio de
        rol, nivel de seguridad o capacidades; concesión y revocación
        de habilitaciones temporales; intentos de acceso no autorizado. 
        \item[-] Cada registro incluye al menos: identificador único del evento,
        marca de tiempo en UTC, identificador del voluntario que realiza
        la acción, tipo de acción, entidad afectada, identificador de la
        entidad afectada, valor anterior (cuando aplique), valor nuevo
        (cuando aplique).
    \end{description}


    \item \textbf{RNF5.2 Inmutabilidad del registro de auditoría} \\
    Los registros de auditoría son inmutables una vez escritos. No debe existir en el sistema ninguna operación que permita modificar
    o eliminar entradas del registro de auditoría.
    \begin{description}
        \item[-] La integridad del registro se protege adicionalmente mediante el
        uso de una tabla con sólo permisos de inserción (no \texttt{UPDATE} ni
        \texttt{DELETE}) a nivel de la base de datos. 
    \end{description}


    \item \textbf{RNF5.3 Consultabilidad del registro} \\
    El registro de auditoría es consultable por usuarios con rol Admin a
    través de una pantalla del sistema.
    \begin{description}
        \item[-] Filtros disponibles: rango de fechas, voluntario, tipo de acción, entidad afectada.
        \item[-] Los resultados se muestran paginados.
        \item[-] Se permite exportación CSV para análisis externo.  
    \end{description}


    \item \textbf{RNF5.4 Trazabilidad de cambios en datos críticos} \\
    Para las entidades especialmente críticas (asientos contables,
    entradas y salidas alimentarias, fichas de amigo en nivel 3) el
    registro de auditoría conserva el valor anterior y el valor nuevo
    de cada campo modificado.
\end{itemize}



\subsection{RNF6 Persistencia y Gestión de Datos}
Este bloque cubre las propiedades del sistema relativas al almacenamiento
y la integridad de la información, así como las copias de seguridad y la
recuperación frente a incidentes.
\begin{itemize}
    \item \textbf{RNF6.1 Sistema de gestión de base de datos} \\
    El sistema utiliza PostgreSQL como base de datos relacional única. La instancia es la 
    provista como servicio gestionado por Render.


    \item \textbf{RNF6.2 Integridad referencial y transaccional} \\
    Las operaciones que afectan a múltiples entidades se ejecutan dentro
    de una transacción atómica.
    \begin{description}
        \item[-] Una entrega de prenda a un amigo (RF5.3) que comporta crear un
        registro de entrega y decrementar el stock se ejecuta como una
        única transacción: o ambos cambios se persisten o ninguno lo hace.
        \item[-] Las invariantes de negocio críticas (un asiento de tipo Amigo
        exige amigo vinculado, una recepción de lavandería exige amigo
        asociado, etc.) se validan tanto en la capa de servicio como
        mediante restricciones declarativas en la base de datos cuando es
        posible.
        \item[-] Las relaciones entre entidades se garantizan mediante claves
        foráneas declaradas a nivel de la base de datos.
    \end{description}


    \item \textbf{RNF6.3 Indexación de columnas de búsqueda frecuente} \\
    Las columnas que participan habitualmente en filtros, ordenaciones
    o cláusulas \texttt{WHERE} deben disponer de los índices apropiados.
    \begin{description}
        \item[-] Como mínimo: \texttt{amigo.dni}` con índice único, \texttt{amigo.nombre}` y
        \texttt{amigo.alias} con índices para búsqueda parcial, fechas en
        asientos contables y entradas/salidas alimentarias, claves foráneas
        en todas las entidades, columnas de filtro frecuente del log de
        auditoría 
    \end{description}


    \item \textbf{RNF6.4 Tipos de datos eficientes} \\
    El diseño del esquema utiliza los tipos de datos más eficientes
    compatibles con los requisitos del dominio.
    \begin{description}
        \item[-] Identificadores numéricos: \texttt{INTEGER} salvo justificación de uso de \texttt{BIGINT}. 
        \item[-] Cadenas de longitud acotada: \texttt{VARCHAR(n)} con \texttt{n} declarado % chktex 36
        explícitamente, no \texttt{TEXT}, salvo en campos de descripción libre. 
        \item[-] Fechas sin componente horario: \texttt{DATE}, no \texttt{TIMESTAMP}.
        \item[-] Booleanos: \texttt{BOOLEAN}.
        \item[-] Importes monetarios: \texttt{NUMERIC(p,s)} con escala fija de dos % chktex 36
        decimales, no \texttt{FLOAT} ni \texttt{DOUBLE}.
    \end{description}
\end{itemize}


\subsection{RNF7 Rendimiento}
Este bloque cubre los criterios cuantitativos de rendimiento del
sistema bajo condiciones de carga esperada.
\begin{itemize}
    \item \textbf{RNF7.1 Tiempos de respuesta} \\
    Los tiempos de respuesta del sistema bajo condiciones normales de
    carga deben mantenerse en niveles que no degraden la experiencia
    de usuario. Estas son métricas orientativas pero deseables en la medida de lo posible:
    \begin{description}
        \item[-] Operaciones simples (consulta de ficha, listado paginado, login): inferior a 500 ms.
        \item[-] Operaciones medias (creación de asiento contable con autocompletado, generación de 
        informe mensual): inferior a 2 segundos.
        \item[-] Operaciones pesadas (importación de datos en lote de hasta 1000
        filas, generación de cuaderno alimentario anual): inferior a 10 segundos. 
    \end{description}


    \item \textbf{RNF7.2 Paginación obligatoria de listados} \\
    Toda consulta que pueda devolver un volumen significativo de
    resultados se entrega paginada.
    \begin{description}
        \item[-] Tamaño de página por defecto: 20 elementos, máximo permitido: 100 elementos. 
        \item[-] Las peticiones que solicitan tamaños superiores se acotan
        silenciosamente al máximo. 
        \item[-] Las consultas devuelven, junto a los datos, los metadatos de
        paginación necesarios para que el cliente pueda navegar
        (página actual, total de páginas, total de registros).
    \end{description}


    \item \textbf{RNF7.3 Optimización de respuestas en backend} \\
    Las respuestas de la API se optimizan para minimizar el tamaño
    transferido y la latencia percibida.
    \begin{description}
        \item[-] Los endpoints utilizan DTOs específicos, no entidades de
        persistencia completas.
        \item[-] Para cada endpoint que devuelve listados, el DTO de listado es
        más ligero que el DTO de detalle. 
    \end{description}


    \item \textbf{RNF7.4 Lazy loading en frontend} \\
    El frontend Angular implementa carga diferida de los módulos no
    imprescindibles para el primer pintado.
    \begin{description}
        \item[-] El bundle inicial contiene únicamente el código necesario para la
        pantalla de login y el armazón de la aplicación autenticada. 
        \item[-] Cada dominio funcional (gestión del amigo, contabilidad,
        trazabilidad alimentaria, etc.) se carga al navegar a él por
        primera vez.
        \item[-] El objetivo es mantener el bundle inicial por debajo de 500 KB comprimidos.
    \end{description}


    \item \textbf{RNF7.5 Generación de PDFs en cliente cuando sea posible} \\
    Los informes en formato PDF que pueden producirse a partir de los
    datos ya cargados en el navegador se generan en el cliente, sin
    involucrar al backend. Esta decisión reduce el ancho de banda saliente del backend y
    distribuye la carga de cómputo al dispositivo del usuario.
    \begin{description}
        \item[-] Se utiliza una librería ligera (jsPDF, pdfmake o equivalente). 
        \item[-] Los informes que requieren acceso a datos no cargados o cálculos
        complejos sí se generan en el backend.
    \end{description}


    \item \textbf{RNF7.6 Operaciones intensivas en segundo plano} \\
    Las operaciones que pueden tardar más de unos pocos segundos se
    ejecutan de forma que no bloquean al usuario cuando sea posible.
    \begin{description}
        \item[-] Las importaciones masivas de datos (RF9.3) se procesan con un
        indicador de progreso visible notificando como ``completado'' cuando finalizan.
    \end{description}
\end{itemize}



\subsection{RNF8 Usabilidad}
Este bloque cubre las propiedades del sistema que determinan la
facilidad y satisfacción con la que los usuarios lo utilizan,
considerando la heterogeneidad del voluntariado real de la asociación.
\begin{itemize}
    \item \textbf{RNF8.1 Diseño responsive} \\
    El sistema funciona correctamente en dispositivos con anchos de
    pantalla desde 320 píxeles (móvil pequeño) hasta 1920 píxeles
    (escritorio amplio).
    \begin{description}
        \item[-] Las operaciones frecuentes en tableta (registrar una ducha,
        entregar una prenda, anotar una entrada alimentaria) están
        optimizadas para uso con una mano. 
        \item[-] Los formularios largos se reorganizan en columna única en
        pantallas estrechas para evitar desplazamientos horizontales o en menús progresivos según convenga.
    \end{description}


    \item \textbf{RNF8.2 Optimización para uso táctil} \\
    En dispositivos táctiles, los elementos interactivos respetan los
    tamaños mínimos recomendados para uso con dedo y acciones potencialmente destructivas (eliminar, anular)
    requieren confirmación explícita para evitar pulsaciones
    accidentales.


    \item \textbf{RNF8.3 Adecuación de la interfaz de usuario} \\
    La interfaz se diseña teniendo en cuenta la heterogeneidad del
    voluntariado, desde perfiles plenamente digitales hasta perfiles
    predominantemente analógicos. Se utilizarán metáforas visuales claras, 
    lenguaje sencillo evitando términos técnicos (CRM, JWT, registro, etc.) 
    y se evitará la sobrecarga de información en formularios o menús.


    \item \textbf{RNF8.4 Mensajes de error útiles} \\
    Los mensajes de error son comprensibles, accionables y orientados
    al usuario, no al desarrollador: debe indicar, cuando es posible, qué hacer para resolver
    el problema.


    \item \textbf{RNF8.5 Consistencia visual y de interacción} \\
    Las pantallas del sistema mantienen consistencia en su disposición,
    nomenclatura, iconografía y patrones de interacción.
    \begin{description}
        \item[-] Los botones de acción primaria y secundaria se distinguen
        visualmente y se ubican consistentemente. 
        \item[-] Los iconos utilizados para una misma acción (editar, eliminar,
        filtrar, exportar) son los mismos en toda la aplicación.
        \item[-] Los formularios siguen el mismo patrón de disposición:
        etiqueta arriba, campo debajo, mensaje de error inmediatamente
        bajo el campo.
    \end{description}


    \item \textbf{RNF8.6 Indicación de progreso y feedback} \\
    El sistema mantiene al usuario informado del estado de las
    operaciones que realiza.
    \begin{description}
        \item[-] Operaciones que toman más de 2 segundos muestran progreso
        cuantitativo cuando es posible (porcentaje, X de Y, etc.). 
        \item[-] Los cambios de estado relevantes (asiento creado, entrega
        registrada, sesión próxima a expirar) se comunican
        proactivamente.
        \item[-] Las operaciones completadas con éxito producen una confirmación
        visual breve y no bloqueante.
    \end{description}


    \item \textbf{RNF8.7 Adecuación lingüística} \\
    La interfaz se ofrece en español de España, con vocabulario y
    giros propios de Andalucía cuando aporten cercanía sin
    comprometer la claridad.
\end{itemize}



\subsection{RNF9 Accesibilidad}
Este bloque cubre las propiedades del sistema que garantizan su
usabilidad por personas con diversidad funcional.
\begin{itemize}
    \item \textbf{RNF9.1 Cumplimiento de WCAG 2.1 AA} \\
    El sistema cumple las pautas de accesibilidad para contenido web
    WCAG 2.2 en su nivel de conformidad AA.\@


    \item \textbf{RNF9.2 Contraste de color} \\
    Los pares texto-fondo cumplen las relaciones de contraste mínimas
    exigidas por WCAG 2.2 AA.\@


    \item \textbf{RNF9.3 Navegación por teclado} \\
    Toda la funcionalidad del sistema es operable mediante teclado, sin
    necesidad de ratón o pantalla táctil.
    \begin{description}
        \item[-] El orden de tabulación sigue una secuencia lógica.
        \item[-] El foco actual es siempre visible.
        \item[-] Las acciones primarias en formularios son ejecutables con la tecla
        \texttt{Enter} desde cualquier campo del formulario.
    \end{description}


    \item \textbf{RNF9.4 Independencia de la interacción del color} \\
    La información transmitida visualmente mediante color se duplica
    mediante otros medios (texto, iconos, patrones).


    \item \textbf{RNF9.5 Adaptabilidad a preferencias del usuario} \\
    La interfaz respeta las preferencias del sistema operativo del
    usuario cuando es razonable hacerlo.
    \begin{description}
        \item[-] Modo oscuro si el sistema operativo lo solicita. 
        \item[-] Tamaño de fuente ajustable según preferencias del sistema operativo.
    \end{description}
\end{itemize}


\subsection{RNF10 Operación y Mantenimiento}
Este bloque cubre las propiedades del sistema relativas a su operación
continuada en producción.
\begin{itemize}
    \item \textbf{RNF10.1 Monitorización de salud} \\
    El sistema expone endpoints de comprobación de salud que permiten
    a Render y a herramientas externas verificar su estado.
    \begin{description}
        \item[-] Endpoint ligero \texttt{/actuator/health} que responde en menos de 100
        ms y no consulta la base de datos, utilizado como health check por la plataforma. 
    \end{description}


    \item \textbf{RNF10.2 Pruebas automatizadas} \\
    El proyecto dispone de un conjunto de pruebas automatizadas que
    verifican el comportamiento esperado del sistema.
    \begin{itemize}
        \item \textbf{RNF10.2.1 Pruebas unitarias} \\
        Las clases de la capa de servicio y los componentes con lógica
        significativa del frontend disponen de pruebas unitarias.
        \begin{description}
            \item[-] Cobertura mínima objetivo: 80\% de líneas para la capa de servicio del backend.
            \item[-] Las pruebas unitarias se ejecutan en cada commit y deben pasar
            para que el push pueda realizarse un merge a las ramas principales. 
        \end{description}
        \item \textbf{RNF10.2.2 Pruebas de integración} \\
        Los flujos críticos del sistema disponen de pruebas de integración
        que verifican el comportamiento extremo a extremo en backend.
        \begin{description}
            \item[-] Flujos cubiertos: autenticación y autorización, creación de
            asiento contable, ciclo de entrega de prenda, registro de entrada
            alimentaria, importación de datos en lote. 
            \item[-] Las pruebas de integración utilizan una base de datos PostgreSQL
            desplegada en contenedor Docker, no mocks.
        \end{description}
        \item \textbf{RNF10.2.3 Pruebas de control de acceso} \\
        El sistema dispone de pruebas específicas que verifican la
        aplicación correcta del control de acceso (RNF3).
        \begin{description}
            \item[-] Cada endpoint dispone de al menos una prueba que verifica que
            un usuario sin las capacidades necesarias recibe respuesta 403. 
        \end{description}
    \end{itemize}


    \item \textbf{RNF10.3 Documentación operativa} \\
    El proyecto mantiene documentación viva sobre su operación.
    \begin{description}
        \item[-] README del repositorio con instrucciones de arranque local,
        ejecución de pruebas, despliegue. 
        \item[-] Documentación de incidentes pasados, sus causas y resoluciones,
        como base de conocimiento.
    \end{description}
\end{itemize}


\subsection{RNF11 Cumplimiento Legal y Normativo}
Este bloque cubre las obligaciones legales que el sistema debe
satisfacer por la naturaleza de los datos y la actividad que gestiona.\
\begin{itemize}
    \item \textbf{RNF11.1 Cumplimiento del Reglamento General de Protección de Datos} \\
    El tratamiento de datos personales realizado por el sistema cumple
    las exigencias del Reglamento (UE) 2016/679 (RGPD) y la Ley
    Orgánica 3/2018 de Protección de Datos Personales y Garantía de los
    Derechos Digitales siguiendo el \textit{principio de limitación de la finalidad}: los datos recogidos para
    una finalidad específica no se utilizan para finalidades incompatibles.


    \item \textbf{RNF11.2 Cumplimiento de la Ley 1/2025 de prevención de pérdidas y desperdicio alimentario} \\
    El módulo de trazabilidad alimentaria (RF6) implementa las exigencias documentales de la Ley 1/2025.
\end{itemize}